% Options for packages loaded elsewhere
\PassOptionsToPackage{unicode}{hyperref}
\PassOptionsToPackage{hyphens}{url}
%
\documentclass[
  10pt,
]{article}
\usepackage{amsmath,amssymb}
\usepackage{iftex}
\ifPDFTeX
  \usepackage[T1]{fontenc}
  \usepackage[utf8]{inputenc}
  \usepackage{textcomp} % provide euro and other symbols
\else % if luatex or xetex
  \usepackage{unicode-math} % this also loads fontspec
  \defaultfontfeatures{Scale=MatchLowercase}
  \defaultfontfeatures[\rmfamily]{Ligatures=TeX,Scale=1}
\fi
\usepackage{lmodern}
\ifPDFTeX\else
  % xetex/luatex font selection
\fi
% Use upquote if available, for straight quotes in verbatim environments
\IfFileExists{upquote.sty}{\usepackage{upquote}}{}
\IfFileExists{microtype.sty}{% use microtype if available
  \usepackage[]{microtype}
  \UseMicrotypeSet[protrusion]{basicmath} % disable protrusion for tt fonts
}{}
\makeatletter
\@ifundefined{KOMAClassName}{% if non-KOMA class
  \IfFileExists{parskip.sty}{%
    \usepackage{parskip}
  }{% else
    \setlength{\parindent}{0pt}
    \setlength{\parskip}{6pt plus 2pt minus 1pt}}
}{% if KOMA class
  \KOMAoptions{parskip=half}}
\makeatother
\usepackage{xcolor}
\usepackage[margin=0.85in]{geometry}
\setlength{\emergencystretch}{3em} % prevent overfull lines
\providecommand{\tightlist}{%
  \setlength{\itemsep}{0pt}\setlength{\parskip}{0pt}}
\setcounter{secnumdepth}{-\maxdimen} % remove section numbering
\usepackage{amsmath,amssymb,amsthm}
\usepackage{fvextra}
\DefineVerbatimEnvironment{verbatim}{Verbatim}{breaklines=true,breakanywhere=true,fontsize=\small}
\ifLuaTeX
  \usepackage{selnolig}  % disable illegal ligatures
\fi
\usepackage{bookmark}
\IfFileExists{xurl.sty}{\usepackage{xurl}}{} % add URL line breaks if available
\urlstyle{same}
\hypersetup{
  pdftitle={Relative energy and the entire Monge--Ampere operator for big rational classes},
  pdfauthor={Reconstructed proof and independent scope review},
  hidelinks,
  pdfcreator={LaTeX via pandoc}}

\title{Relative energy and the entire Monge--Ampere operator for big
rational classes}
\author{Reconstructed proof and independent scope review}
\date{6 October 2026}

\DeclareUnicodeCharacter{2212}{\ensuremath{-}}
\DeclareUnicodeCharacter{2295}{\ensuremath{\oplus}}
\DeclareUnicodeCharacter{2264}{\ensuremath{\le}}
\DeclareUnicodeCharacter{03B1}{\ensuremath{\alpha}}
\DeclareUnicodeCharacter{2211}{\ensuremath{\sum}}
\DeclareUnicodeCharacter{27E8}{\ensuremath{\langle}}
\DeclareUnicodeCharacter{27E9}{\ensuremath{\rangle}}
\DeclareUnicodeCharacter{00B7}{\ensuremath{\cdot}}
\DeclareUnicodeCharacter{2265}{\ensuremath{\ge}}
\DeclareUnicodeCharacter{03BC}{\ensuremath{\mu}}
\DeclareUnicodeCharacter{03B4}{\ensuremath{\delta}}
\DeclareUnicodeCharacter{222B}{\ensuremath{\int}}
\DeclareUnicodeCharacter{2032}{\ensuremath{\prime}}
\DeclareUnicodeCharacter{2192}{\ensuremath{\to}}
\DeclareUnicodeCharacter{03BB}{\ensuremath{\lambda}}
\DeclareUnicodeCharacter{03A9}{\ensuremath{\Omega}}
\DeclareUnicodeCharacter{03C6}{\ensuremath{\varphi}}
\DeclareUnicodeCharacter{2193}{\ensuremath{\downarrow}}
\DeclareUnicodeCharacter{03B5}{\ensuremath{\epsilon}}
\DeclareUnicodeCharacter{221E}{\ensuremath{\infty}}
\DeclareUnicodeCharacter{03B3}{\ensuremath{\gamma}}
\DeclareUnicodeCharacter{03BA}{\ensuremath{\kappa}}
\DeclareUnicodeCharacter{00B1}{\ensuremath{\pm}}
\DeclareUnicodeCharacter{03B8}{\ensuremath{\theta}}
\DeclareUnicodeCharacter{0394}{\ensuremath{\Delta}}
\DeclareUnicodeCharacter{2208}{\ensuremath{\in}}
\DeclareUnicodeCharacter{2191}{\ensuremath{\uparrow}}
\DeclareUnicodeCharacter{2209}{\ensuremath{\notin}}
\DeclareUnicodeCharacter{2113}{\ensuremath{\ell}}
\DeclareUnicodeCharacter{03C3}{\ensuremath{\sigma}}
\begin{document}
\maketitle

\section{Abstract and exact theorem}\label{abstract-and-exact-theorem}

Let \(X\) be a smooth integral projective variety of dimension \(n\ge1\)
over an algebraically closed trivially valued field of characteristic
zero and arbitrary cardinality. Let \(L\) be any actual big rational
line bundle. No nefness, finite generation or finite Zariski
decomposition is assumed. The potentials are the literal genuine
numerical psh functions on the whole original \(X^{\mathrm{an}}\),
defined by single decreasing nets of actual section metrics with
numerical class tending to \(c_1(L)\).

Put \(V_L=P_L(0)\) and \(v=\operatorname{vol}_X(L)>0\). This paper
constructs a target-independent section/volume energy on the entire
genuine cone, normalized by \(E_L(V_L)=0\), with \(E_L(u+c)=E_L(u)+cv\).
The intrinsic relative domain is

\[\mathcal E^1(L,V_L)=\{u\in\operatorname{PSH}(L):E_L(u)>-\infty\}.\]

Its energy is characterized by the genuine reference floors
\(\max(u,V_L-k)\), is concave and upper semicontinuous in the original
divisorial topology, and has compact normalized energy-cost sublevels.
On that entire domain the construction gives a positive Radon metric
operator \(M_L(u)\) of raw mass \(v\), agreeing with the actual
continuous-obstacle section/volume operator. It is independent of
decreasing model-majorant nets, including noncofinal nets with the same
all-point endpoint. It satisfies both continuous-direction outer
derivative signs:

\[\left.\frac{d}{dt}\right|_{t=0}E_L(P_L(u+th))
   =\int h\,dM_L(u),\qquad h\in C(X^{\mathrm{an}}).\]

Every envelope here is an actual original genuine function. The operator
is weakly continuous along arbitrary decreasing genuine nets with
finite-energy endpoint, and under original divisorial convergence plus
energy convergence. All measures and continuous observables remain on
the original analytic carrier.

Xean independent review128, using native Codex gpt-6.1-sol with xhigh
reasoning, accepted precisely this theorem. It did not review mixed
products, integrable weighted limits, pole-nullness, comparison,
domination or an all-target inverse equation. Consequently original
P16.5 remains open. The reviewed theorem is a supporting result for its
full equation. The five complete proof units below retain their frozen
mathematical bodies; earlier draft labels and historical gap sentences
indicate when those units were written. This abstract and the complete
independent report at the end specify the accepted scope.

\section{0. Inputs, complete section limits and the sole ample
grant}\label{inputs-complete-section-limits-and-the-sole-ample-grant}

The exact actual-ample G0 theorem is the supplied equation premise. It
is invoked only inside an actual ample class for genuine closure,
divisorial compactness and all-point determination. Its ample
intersection mass, energy, mixed products and inverse equation are never
relabeled as nonample statements. The exact full source is the accepted
P16.1 round-92 mathematical argument, bound in the source ledger. None
of its ample finite-energy memberships is assumed for an unbounded
nonnef potential.

The general envelope/reference premise is Boucksom--Jonsson's addendum,
arXiv:2206.07183v2, Theorem A and Corollaries 2.1--2.2. Ordinary
projective positive products and the volume differential are used in the
smooth characteristic-zero domain of Boucksom--Favre--Jonsson,
arXiv:math/0608260v2. These are published mathematical inputs; their
precise uses occur in the proof units and review report.

In the lattice formulas one needs a complete section limit, rather than
just the limsup definition of volume. For a big Cartier divisor on a
projective variety the required statement is

\[\lim_{m\to\infty}\frac{d!\,h^0(Y,mA)}{m^d}=\operatorname{vol}_Y(A).\]

The complete-series argument in Lazarsfeld--Mustata, \emph{Convex bodies
associated to linear series}, Sections 2.1--2.2, Proposition 2.1, Lemma
2.2 and Theorem 2.3, is the published input for this limit. Its original
uncountable-field convention is removed at this particular application
by choosing an uncountable algebraically closed extension \(K/k\). Flat
cohomology base change preserves each actual section dimension
\(h^0(Y,mA)\) exactly. Thus it preserves the complete numerical
sequence, its limsup and the obtained limit. No countability restriction
is added to the original field. Clearing denominators gives the
corresponding limit through divisible degrees for actual rational
bundles. This supplementary application supplies the full section-limit
input used throughout the frozen units, without pretending that BFJ's
limsup definition alone proves convergence.

Smooth equivariant common refinements are obtained by resolving the
closure of the generic graph and its central boundary, using functorial
characteristic-zero resolution. Normality and the height-one
characterization of regularity justify the section identities and
exceptional coefficient vanishing. The precise imported propositions,
original domains and source links are recorded in the independent report
and source ledger. The proof is relative to these foundations and the
exact ample grant; it does not claim a recursive reconstruction of every
external foundation.

\section{I. Genuine minimal reference and the volume
obstruction}\label{i.-genuine-minimal-reference-and-the-volume-obstruction}

\textbf{Status: supplementary research note, independently unreviewed.}
The original P16.5 remains open. Frozen native124 and native125 are
unchanged. This note updates the remaining research obligations in
author123: the general-reference existence gate has a published primary
input, while the relative operator and variational equation still
require construction.

\subsubsection{1. Exact imported primary
result}\label{exact-imported-primary-result}

Boucksom--Jonsson, \href{https://arxiv.org/pdf/2206.07183v2}{Addendum to
Global pluripotential theory over a trivially valued field}, version2
dated September2023, proves the envelope property for every numerical
class on smooth projective X over algebraically closed trivially valued
k when char(k)=0. This retains arbitrary cardinality. Corollaries2.1
and2.2 provide the genuine reference V\_theta=P\_theta(0), with generic
value zero, and compactness of the normalized genuine cone. These are
external results, not an additional equation grant.

The complete pertinent proof units were opened: Lemmas1.3--1.7, rendered
lines80--257; TheoremA's proof and Corollaries2.1--2.3, lines324--405.
The explicit multiplier-ideal approximants in the big rational case form
a decreasing divisibility net in actual L+A/m, so the required original
single-net membership must be checked directly rather than delegated to
the disputed abstract flattening lemma. No full source transcript is
published here.

\subsubsection{2. Minimality and the reference actually
retained}\label{minimality-and-the-reference-actually-retained}

For an actual big rational L and theta=c1(L), choose any genuine
theta-potential u. A common larger actual ample cone gives
u(v\_triv)=sup u. Normalize that supremum to zero. Then u is a
competitor for P\_theta(0), so \[
 u\le V_\theta\le0.
\] Consequently V\_theta is a genuine potential with minimal
singularities. It is intrinsic to the original numerical cone and
canonical metric convention; an arbitrary effective divisor in an
ample-plus-effective decomposition is not substituted for it.

For any usc obstacle h with a genuine minorant, the external envelope
property ensures the regularized supremum of genuine competitors is
genuine; because h is usc, it remains below h and is itself a
competitor. Thus P\_theta(h) is the genuine obstacle envelope. In
particular max(u,V\_theta-k) belongs to the original cone. These
reference floors are genuine functions. That fact does not alone
construct a positive relative operator, prove a finite-energy
definition, or identify the operator's mass with volume.

\subsubsection{3. Why the nef raw-intersection recipe fails for a big
nonnef
bundle}\label{why-the-nef-raw-intersection-recipe-fails-for-a-big-nonnef-bundle}

Let X be the blowup of P2\_k at one point, H the pullback of O(1), and E
the exceptional divisor. Take the actual integral bundle L=2H+E. The
intersection numbers H\^{}2=1, H.E=0 and E\^{}2=-1 give \[
 L.E=-1,\qquad L^2=3.
\] Thus L is nonnef. For each m\textgreater=1 and j=1,\ldots,m, \[
 (2mH+jE)|_E\simeq O_{P^1}(-j).
\] The Cartier restriction sequence and the absence of sections of O(-j)
show that multiplication by the canonical E-section identifies \[
 H^0(X,2mH+(j-1)E)\simeq H^0(X,2mH+jE).
\] Iterating gives H\textsuperscript{0(X,mL)=s\_E}m H\^{}0(X,2mH).
Proper birational normal descent identifies the latter section space
with H\^{}0(P2,O(2m)). Hence \[
 h^0(X,mL)=\binom{2m+2}{2},\qquad \operatorname{vol}(L)=4>0.
\] This is a big nonnef example on the retained smooth projective
original carrier for which volume differs from the ordinary top
intersection. Replacing the mass4 by L\^{}2=3 would produce the wrong
equation.

The nef proof's canonical zero is also unavailable here: a bounded
endpoint of the genuine numerical FS definition would force every
defining stage to have no common zeros. Those sections generate their
bundle, making each stage's numerical class nef. Its numerical limit
would then be nef, contradicting L.E=-1. In particular zero is not a
genuine L-potential.

\subsubsection{4. Remaining construction}\label{remaining-construction}

The next task is to construct the original relative finite-energy
operator using V\_theta, with mass vol(L), and prove its weighted
limits, comparison, outer differentiation and actual attainment.
Fixed-plus-movable intersection approximants on modifications need an
original-carrier independence and descent proof before they can define
that operator. Neither normalized cone compactness nor genuine reference
floors supply those conclusions by themselves. The separate nonbig
additional-data and uniqueness problem also remains open.

\section{II. Actual lattices, volume energy and atomic
mass}\label{ii.-actual-lattices-volume-energy-and-atomic-mass}

This is an unreviewed supporting argument for original P16.5. It
constructs a target-independent functional on vertical model obstacles
and its actual finite atomic derivative. It does not yet identify the
functional on the entire genuine relative finite-energy cone, prove
descent through equal psh envelopes, or solve the original equation.
Original P16.5 remains open.

\subsubsection{1. Exact objects and imported algebraic
premises}\label{exact-objects-and-imported-algebraic-premises}

Let X be smooth integral projective of dimension n over an algebraically
closed field k of characteristic zero, trivially valued, and let L be an
actual big rational line bundle. Put V=vol\_X(L)\textgreater0 and d=n+1.
No nefness, finite generation, or birational Zariski decomposition is
assumed.

A model here is a smooth projective variety Y with a projective
birational morphism rho:Y -\textgreater{} X times P\^{}1, equivariant
for the multiplicative group and an isomorphism away from the fiber over
0. Let pi be the map to P\^{}1, F=pi\textsuperscript{\emph{(0), and D a
rational Cartier divisor supported on F. Set B\_D=rho\^{}}p\_X}*L+D and
alpha\_D(c)=B\_D+cF. All tensor identifications use these actual line
bundles. Degrees m below are divisible by a fixed denominator.

The algebraic background used explicitly is: full section asymptotics of
a big rational bundle have a limit; volume is birationally invariant and
monotone under adding effective divisors; volume is continuous; and on
the big cone its differential is d times the positive intersection
product. The volume root of degree 1/d is concave. These are imported
algebraic premises, not additional grants of a non-Archimedean equation.

Primary evidence for the differential is Boucksom--Favre--Jonsson,
arXiv:math/0608260v2, Theorem A and its complete proof in Section 3.2,
pp.16--17. Their construction and continuity proof of positive products
are in Section 2.2--2.3; Theorem 2.16 states root concavity. The
section-asymptotic/Fujita and Hodge-index foundations are not
recursively proved in this note. Thus this note is not an independently
accepted complete dependency chain.

\subsubsection{2. Exact lattice calculation and eventual affine
volume}\label{exact-lattice-calculation-and-eventual-affine-volume}

Write F=sum b\_i E\_i, with the actual positive central multiplicities.
Since D is supported on F, choose rational a,b with aF \textless= D
\textless= bF coefficientwise. For each divisible m, let
R\_m=H\^{}0(X,mL), of dimension r\_m, and consider the lattice

\begin{verbatim}
M_m=H^0(pi^(-1)(A^1),mB_D)
    inside R_m tensor k[t,t^(-1)].
\end{verbatim}

Properness makes pi\_\emph{O\_Y(mB\_D) coherent. It is torsion free over
the nonsingular base: multiplication by a nonzero base function is
injective on sections of a line bundle on the integral Y. It is
therefore locally free. Away from 0, rho identifies its restriction with
R\_m tensor O, and the chosen trivialization at infinity is the one
coming from the actual product bundle. Birational normality gives
rho\_}O\_Y=O and hence the same sections for the pullback product
bundle. The divisor inequalities give actual lattice inclusions

\begin{verbatim}
t^(-ma) R_m[t] <= M_m <= t^(-mb) R_m[t].
\end{verbatim}

The equivariance makes M\_m graded for the t degree in this fixed
product trivialization. A finite filtration of the vector space R\_m is
obtained by taking the coefficients of its homogeneous pieces. Choose a
basis adapted to this filtration. It gives integers w\_m,j such that

\begin{verbatim}
M_m = direct sum_j k[t] t^(-w_m,j) s_m,j,
ma <= w_m,j <= mb.
\end{verbatim}

Equivalently, glue this lattice to the fixed product lattice near
infinity: the resulting vector bundle on P\^{}1 is the direct sum of
O(w\_m,j). This also proves the splitting needed here without an
assertion that the entire section algebra is finitely generated.
Projection formula and the definition of direct image yield

\begin{verbatim}
h^0(Y,m alpha_D(c)) = sum_j max(w_m,j+mc+1,0).
\end{verbatim}

For rational c\textgreater-a and sufficiently divisible m every summand
is positive. Consequently the exact formula is

\begin{verbatim}
h^0(Y,m alpha_D(c)) = (mc+1)r_m + sum_j w_m,j.             (2.1)
\end{verbatim}

It also shows that alpha\_D(c) is big: its volume is at least
d(c+a)V\textgreater0. For any rational h with c+h\textgreater-a,
subtracting the two exact formulas and using the full-section asymptotic
r\_m/m\^{}n -\textgreater{} V/n! gives

\begin{verbatim}
vol_Y(alpha_D(c+h))-vol_Y(alpha_D(c)) = d h V.             (2.2)
\end{verbatim}

Both big bundles have genuine asymptotic volume limits under the stated
algebraic premise; there is no exchange of two uncontrolled limsups.
Continuity extends (2.2) to real c,h in this half-line. In particular,

\begin{verbatim}
vol_Y(alpha_D(c)) = d c V + A_Y(D)
\end{verbatim}

for all sufficiently large c.~The intercept is finite and is calculated
from actual section lattices, including horizontal base loci of L. The
argument never assumes that B\_D is relatively generated or that its
base locus is confined to the central fiber.

\subsubsection{3. A well-defined algebraic energy of
obstacles}\label{a-well-defined-algebraic-energy-of-obstacles}

For c large enough for both terms define

\begin{verbatim}
E_L^alg(D) = [vol_Y(alpha_D(c))-vol_Y(alpha_0(c))]/d.
\end{verbatim}

Equation (2.2) proves independence of c.~On the product the exact
Kunneth section count gives vol(alpha\_0(c))=dcV; birational invariance
gives the same equality on Y. Hence E\_L\^{}alg(D)=A\_Y(D)/d and
E\_L\^{}alg(0)=0.

If f:Y' -\textgreater{} Y is an equivariant birational morphism of such
smooth models, replace D by f\^{}\emph{D and F by f\^{}}F. Birational
invariance of both volumes proves equality of energies on the two
models. Any two permitted models have a common smooth equivariant
resolution; thus the value is independent of the chosen representation
of a Cartier vertical divisor on all models.

For real s, coefficientwise effective vertical H, and rational C
satisfying -CF \textless= D-D' \textless= CF, the following hold:

\begin{verbatim}
E_L^alg(D+sF)=E_L^alg(D)+sV,
E_L^alg(D+H)>=E_L^alg(D),
|E_L^alg(D)-E_L^alg(D')|<=CV.                             (3.1)
\end{verbatim}

The first uses (2.2); the second is volume monotonicity; the third
follows by sandwiching with the first two. Irrational coefficients are
obtained by continuity on the fixed finite-dimensional vertical-divisor
space. No ordering of residual section metrics is used.

The functional is concave. Fix D\_0,D\_1 and s in {[}0,1{]}, set
D\_s=(1-s)D\_0+sD\_1, and choose a common sufficiently large
c.~Volume-root concavity gives

\begin{verbatim}
[dcV+dE(D_s)]^(1/d)
  >= (1-s)[dcV+dE(D_0)]^(1/d)+s[dcV+dE(D_1)]^(1/d).
\end{verbatim}

Divide by (dcV)\^{}(1/d), expand (1+x/c)\^{}(1/d) at zero, multiply by c
and let c tend to infinity. Every intercept is fixed and finite. The
result is E(D\_s)\textgreater=(1-s)E(D\_0)+sE(D\_1). This passes from a
root-concave (n+1)-dimensional volume to an ordinary concave vertical
energy and is not a claim that ordinary volume itself is concave.

\subsubsection{4. Actual positive central coefficients and their
independence}\label{actual-positive-central-coefficients-and-their-independence}

For c large let

\begin{verbatim}
q_i(D,c)=<alpha_D(c)^n> dot E_i.
\end{verbatim}

This number is nonnegative. Indeed for r\textgreater=0, adding rE\_i
increases volume, and its right differential at zero is d q\_i by the
imported C1 theorem. Since F=sum b\_i E\_i, linearity of this
differential and (2.2) give

\begin{verbatim}
sum_i b_i q_i(D,c) = <alpha_D(c)^n> dot F = V.            (4.1)
\end{verbatim}

Each individual coefficient is independent of sufficiently large c.~For
t in a fixed neighborhood of zero, the coefficient bounds
aF\textless=D+tE\_i\textless=bF can be chosen uniformly in t by
enlarging a,b. Thus

\begin{verbatim}
vol(alpha_D(c)+tE_i)=dcV+A_Y(D+tE_i)
\end{verbatim}

for all c in one common half-line. Differentiating in t at zero proves
that q\_i(D,c)=q\_i(D), independent of c.~This argument is necessary:
independence of the total mass alone would not prove independence of
each summand.

Let x\_i be the semivaluation on X obtained by restricting
b\_i\^{}(-1)ord\_Ei through the Gauss extension; repetitions, if any,
are combined as measures. Define the actual Radon measure

\begin{verbatim}
mu_D^alg = sum_i b_i q_i(D) delta_xi.                     (4.2)
\end{verbatim}

It is positive, target independent, and has raw mass vol(L). It is not
declared to be the Monge--Ampere measure of an arbitrary genuine psh
function.

This atomic measure is invariant under replacing Y by a higher smooth
model and D by its pullback. To see this, if G is an effective
exceptional divisor of f:Y' -\textgreater{} Y, normality of Y gives
f\_*O\_Y'(kG)=O\_Y for all k\textgreater=0: a rational function with
poles only over a locus of codimension at least two is regular on normal
Y. Projection formula implies

\begin{verbatim}
vol_Y'(f^*alpha+rG)=vol_Y(alpha), r>=0.
\end{verbatim}

The C1 derivative at r=0 yields
\textless f\^{}\emph{alpha\^{}n\textgreater{} dot G=0. If E\_i' is a
strict transform, f\^{}}E\_i-E\_i' is effective exceptional, so the same
derivative and pullback invariance of volume give

\begin{verbatim}
<f^*alpha^n> dot E_i'=<alpha^n> dot E_i.
\end{verbatim}

Strict transforms retain their central multiplicities and restricted
semivaluations; all new exceptional components have zero coefficient.
Hence (4.2) gives the same measure after refinement. This is a proof of
model invariance, not an assumption that every new central atom has
positive weight.

For another vertical Cartier divisor H=sum h\_i E\_i on the same model,
the exact two-sided derivative is

\begin{verbatim}
derivative at t=0 of E_L^alg(D+tH)
    = <alpha_D(c)^n> dot H
    = sum_i h_i q_i(D).                                 (4.3)
\end{verbatim}

Writing the usual vertical model function as f\_H, with
f\_H(x\_i)=h\_i/b\_i and constants f\_(sF)=s, this is integral f\_H d
mu\_D\^{}alg. This convention is the additive metric of O(H) using its
section on the punctured family; the effective-divisor value is
positive. It is distinct from the function
log\textbar ideal\textbar=-v(ideal).

If the full original-carrier uniform density and representation theorem
for rational vertical model functions is supplied, (3.1) uniquely
extends E\_L\^{}alg to a monotone concave V-Lipschitz functional on
C\textsuperscript{0(X}an). At every model obstacle D, (4.3) then gives
its two-sided derivative in every original continuous direction:
approximate f uniformly by a model f\_H, bound the difference of
difference quotients by
V\textbar\textbar f-f\_H\textbar\textbar\_infinity, and bound the
difference of integrals by the same quantity. First let t tend to zero
for this fixed H, then improve the approximation. This last conclusion
explicitly retains the density/representation theorem as a further
premise; it does not supply differentiability at every continuous
obstacle or every relative finite-energy base.

\subsubsection{5. The remaining original
obligations}\label{the-remaining-original-obligations}

The construction has the correct raw volume mass even for big nonnef L
in any dimension, and it uses its actual line bundle and actual section
lattices. In the example L=2H+E on Bl\_p P\^{}2 from own141 its mass is
4, never the ordinary top intersection 3.

However, to obtain the original metric operator one must still prove at
least the following.

\begin{enumerate}
\def\labelenumi{\arabic{enumi}.}
\tightlist
\item
  The relation between this obstacle energy and the genuine psh envelope
  P\_L(f\_D) from own141. In particular, equality of two envelopes must
  imply equality of the energy and atomic measure. Birational invariance
  proved above is not this envelope-fiber descent.
\item
  A relative energy on the entire original cone with the canonical
  singular reference V\_L, and agreement of its bounded-relative stages
  with the actual model construction. A continuous obstacle need not
  produce a bounded original psh function when L is nonnef.
\item
  Target-independent mixed products, weighted limits, exact relative
  floors, comparison and domination, with pole sets and all original
  continuous observables retained.
\item
  The intrinsic target criterion, weak energy upper semicontinuity,
  actual attainment, and all-point uniqueness modulo constants. Scalar
  concavity and the model derivatives alone do not prove these
  assertions.
\item
  Any additional nonbig range separately. Nothing here normalizes a
  zero-volume measure.
\end{enumerate}

The indispensable imported algebraic asymptotic and differentiability
inputs, and the optional model-function density input, require their own
complete application checks before independent acceptance. The newer
entropy approximation theorem of Trusiani, arXiv:2605.30063v2, Appendix
A.2, explicitly keeps an ample generic line bundle over k((t)); it
supplies no missing general-big equation in this argument. Its
projective Fujita regimes are already established subcases in the
original P19.2 dossier and do not close that residual hybrid problem.

\section{III. Envelope descent and continuous-obstacle
derivatives}\label{iii.-envelope-descent-and-continuous-obstacle-derivatives}

This continues unreviewed own143. It removes the envelope-fiber
obstruction for continuous obstacles and constructs an actual
target-independent operator on their genuine psh envelopes. Its
supporting algebraic premises and the published general envelope
property must still be independently reconstructed. It does not yet
cover the entire relative E1 cone or establish the original variational
equation. Original P16.5 remains open.

\subsubsection{1. Original-carrier conventions and precise
inputs}\label{original-carrier-conventions-and-precise-inputs}

Keep exactly X,k,L,n,V,Y,F,D and the raw algebraic energy of own143. Let
q\_s=m\^{}(-1)log\textbar s\textbar{} for a nonzero actual section s of
mL, with the canonical trivial metric. It is a generically finite
Fubini--Study function in the actual L class: the constant net
consisting of this function meets the defining genuine-psh
approximation. Its values are \textless=0 and its generic value is zero,
including the -infinity convention at section poles.

Use the general numerical envelope property for the smooth
characteristic-zero domain, as in own141. For any continuous real
obstacle f, P\_L(f) is the genuine usc supremum of L-psh minorants, and
P\_L(f)\textless=f.~It exists and is generically finite because
V\_L=P\_L(0) is genuine and f is bounded. Consequently

\begin{verbatim}
V_L-||f||_infinity <= P_L(f) <= V_L+||f||_infinity.
\end{verbatim}

Use also the full original-carrier density and vertical representation
of rational model functions. Primary selected units of
Boucksom--Jonsson, arXiv:1801.08229v4, Section 2.1--2.2, Theorems 2.2
and 2.7 with their complete proofs, have been read in the local source,
literal LF943--1003 and LF1026--1072. Single-section genuine membership
uses Definition 4.1, LF2174--2203. These are precise applications in the
same smooth projective trivially valued setting; arbitrary original C0
tests and semivaluation supports remain present.

The general envelope property is an external published input from
arXiv:2206.07183v2. It is not a grant of the conclusion proved below.
The algebraic positive-product premises, their monotonicity and nef
pairings are as in own143 and the complete pertinent proof units of
arXiv:math/0608260v2, Section 2.2--2.3 and Section 3.2. Recursive
Fujita, Hodge-index and resolution foundations are still explicit
background obligations.

\subsubsection{2. Genuine envelopes determine every actual section
filtration}\label{genuine-envelopes-determine-every-actual-section-filtration}

For an original continuous f and nonzero s in R\_m=H\^{}0(X,mL), set

\begin{verbatim}
W_m(f;s)=inf over x in X^an of [mf(x)-log|s|(x)].
\end{verbatim}

At poles of s the bracket is +infinity. Since
log\textbar s\textbar\textless=0 and its generic value is zero,
W\_m(f;s) is finite and lies in
{[}-m\textbar\textbar f\textbar\textbar\_infinity,m\textbar\textbar f\textbar\textbar\_infinity{]}.
For every real lambda,

\begin{verbatim}
W_m(f;s)>=lambda
  iff q_s+lambda/m <= f
  iff q_s+lambda/m <= P_L(f).                            (2.1)
\end{verbatim}

The last equivalence is exact on the whole carrier. Its forward
implication follows because q\_s+lambda/m is itself a genuine psh
minorant; the reverse follows from P\_L(f)\textless=f.~No density
assertion about all psh functions being generated by sections of this
fixed L is needed.

For a vertical model obstacle f\_D, the condition for an integer lambda
is equivalently that s t\^{}(-lambda) is a regular section of mB\_D on Y
over A\^{}1. The only potentially negative divisor coefficients of this
rational section are central: s is already regular on the generic
product, and horizontal pullback zero divisors are effective. At E\_i
the condition is

\begin{verbatim}
ord_Ei(s)+m ord_Ei(D)-lambda b_i >= 0,
\end{verbatim}

which is precisely (2.1) evaluated at x\_i. Conversely, normality turns
nonnegative coefficients at all prime divisors into regularity. A
regular section has nonnegative value for every centered semivaluation,
so the condition also holds for every point of X\^{}an. Thus this is an
equivalence, not merely an inequality checked at the central atoms.

Therefore the homogeneous lattice of own143, including each of its
filtered coefficient subspaces, is determined by the whole genuine
envelope P\_L(f\_D). If two model obstacles have the same envelope, all
their actual filtered section spaces are identical in the fixed product
trivialization, in every divisible degree. Their raw energies are equal.

More generally P\_L(f\_D)\textless=P\_L(f\_H) implies inclusion of the
coefficient subspaces at every integer level, hence ordered weight sums,
and hence

\begin{verbatim}
E_L^alg(D)<=E_L^alg(H).                                  (2.2)
\end{verbatim}

The weight-sum implication can be seen by choosing a common lower and
upper integer bound for the weights and summing dimensions of the
decreasing flags between those bounds. It does not require one basis to
be simultaneously adapted to two flags. Adding a rational constant c
gives

\begin{verbatim}
P_L(f_D)<=P_L(f_H)+c implies E(D)<=E(H)+cV.
\end{verbatim}

\subsubsection{3. Actual continuous-obstacle section
energy}\label{actual-continuous-obstacle-section-energy}

For every real lambda, the set \{0\} union
\{s:W\_m(f;s)\textgreater=lambda\} is a k-vector subspace. Indeed the
non-Archimedean triangle inequality for actual section norms gives
W\_m(f;s+s')\textgreater=min(W\_m(f;s),W\_m(f;s')), and multiplication
by a nonzero scalar of the trivially valued k leaves W\_m unchanged.
This finite-dimensional decreasing flag admits an adapted basis. Its
r\_m weights, with multiplicities, are intrinsic; the weak inequalities
in the definition retain their endpoints. Define

\begin{verbatim}
e_m(f) = n! m^(-(n+1)) sum_j floor(W_m(f;s_m,j)).
\end{verbatim}

For a rational model f\_D these are exactly the integer lattice weights
of own143, so e\_m(f\_D)-\textgreater E(D). If
\textbar\textbar f-g\textbar\textbar\_infinity\textless=epsilon, then
every section's weights differ by at most m epsilon, and flag comparison
gives

\begin{verbatim}
|e_m(f)-e_m(g)| <= (n! r_m/m^n)(epsilon+1/m).             (3.1)
\end{verbatim}

Uniform model approximation of this one continuous f and the
full-section limit n!r\_m/m\^{}n -\textgreater{} V prove that e\_m(f)
has a limit. Explicitly fix a model g within epsilon of f, use
convergence of e\_m(g), and then use (3.1); the upper and lower limits
for e\_m(f) differ by at most 2V epsilon. Let epsilon decrease to zero.
This uses sequences only for approximation of a single continuous
function in the uniform norm, not to replace arbitrary psh nets on a
nonmetrizable carrier.

Write E\_L\^{}obs(f)=lim e\_m(f). It is the unique V-Lipschitz
continuous extension of the model energy. It is finite, monotone,
concave, has E(0)=0 and E(f+c)=E(f)+cV, by own143 and uniform
approximation. Equations (2.1) and flag comparison now prove, for all
continuous f,g,

\begin{verbatim}
P_L(f)<=P_L(g) implies E_L^obs(f)<=E_L^obs(g),
P_L(f)=P_L(g) implies E_L^obs(f)=E_L^obs(g).               (3.2)
\end{verbatim}

Thus the energy of a genuine continuous-obstacle envelope is already
well defined. It is computed from actual sections and canonical trivial
norms before any target measure is chosen.

\subsubsection{4. Uniform two-sided model derivatives on bounded
obstacle
sets}\label{uniform-two-sided-model-derivatives-on-bounded-obstacle-sets}

The following estimate supplies the missing continuity across models.
Fix a rational vertical model direction h=f\_H and C\textgreater0.
Represent H on one fixed smooth model Z. On any common refinement, a
model obstacle f\_D with
\textbar\textbar f\_D\textbar\textbar\_infinity\textless=C has
-CF\textless=D\textless=CF, by the model-function effectivity criterion.
Choose a fixed sufficiently large rational c.~On the Riemann--Zariski
space of Z put

\begin{verbatim}
alpha_- = pullback(L)+(c-C)F,
alpha_+ = pullback(L)+(c+C)F,
alpha_D = pullback(L)+D+cF.
\end{verbatim}

The first is big by the exact lattice calculation of own143; and
alpha\_-\textless=alpha\_D\textless=alpha\_+ in the pseudoeffective
order on every common refinement. Since alpha\_- is big in the fixed
finite-dimensional N1(Z), there exists a\textgreater0 with a alpha\_-
\textgreater= +/-H. For epsilon=a\textbar t\textbar\textless1,

\begin{verbatim}
(1-epsilon)alpha_D <= alpha_D+tH <= (1+epsilon)alpha_D.
\end{verbatim}

Let Q\_D=\textless alpha\_D\^{}n\textgreater. Monotonicity and
homogeneity of positive products yield

\begin{verbatim}
(1-epsilon)^n Q_D <= <(alpha_D+tH)^n>
                       <= (1+epsilon)^n Q_D.            (4.1)
\end{verbatim}

Choose fixed nef Cartier classes T\_1,T\_2 on Z with H=T\_1-T\_2, by
adding a sufficiently ample class. Pairing a pseudoeffective Weil curve
class with a nef Cartier class is nonnegative. Pair (4.1) with each T\_i
and use Q\_D\textless=Q\_+=\textless alpha\_+\^{}n\textgreater. This
gives a finite constant K, depending only on the fixed data C,H,c and
the positive products on Z, such that

\begin{verbatim}
|[<(alpha_D+tH)^n>-Q_D] dot H| <= K|t|.                 (4.2)
\end{verbatim}

No choice of the model carrying D enters K. The possible numerical
traces on higher models do not have to be nef. The relevant pairing is
the nef-Cartier/positive-Weil pairing on all models, not an intersection
of arbitrary Weil divisors.

The C1 algebraic volume differential and the fundamental theorem of
calculus on this fixed one-dimensional big-class segment, followed by
own143, now give

\begin{verbatim}
|E(f_D+th)-E(f_D)-t integral h d mu_D^alg| <= K t^2/2.  (4.3)
\end{verbatim}

Choose c large enough for the common small-t coefficient bound, so all
uses of the intercept formula occur on the same permitted half-line.
This proves a uniform error on bounded original model obstacles. Merely
knowing separate differentiability at each model would not justify the
next passage.

\subsubsection{5. A genuine continuous-obstacle envelope
operator}\label{a-genuine-continuous-obstacle-envelope-operator}

Fix any continuous f and take a net of model obstacles f\_D tending
uniformly to f.~Their actual measures mu\_D\^{}alg are positive of mass
V on the compact original X\^{}an. Their compact weak-star measure space
has cluster points.

For any fixed model direction h, pass (4.3) through uniform convergence
of f\_D and the continuous energy. For any cluster point mu one gets

\begin{verbatim}
|E(f+th)-E(f)-t integral h dmu| <= K t^2/2.
\end{verbatim}

Thus every cluster point has the same h integral, namely the two-sided
derivative of E at f in direction h. Density of model functions in the
full original C0 and the common finite mass make this cluster point
unique. Therefore the entire net mu\_D\^{}alg converges weakly to one
positive Radon measure mu\_f of mass V, independently of the
approximating net. For every original continuous direction h, uniform
model approximation of h and the V-Lipschitz energy bound yield

\begin{verbatim}
derivative at t=0 of E_L^obs(f+th) = integral h dmu_f.    (5.1)
\end{verbatim}

This is a proved two-sided Gâteaux derivative, not a claim of Fréchet
differentiability in the uniform norm. The same bounded-set argument
proves continuity of f -\textgreater{} mu\_f for the uniform topology on
obstacles and the full original weak measure topology.

The measure descends through the genuine envelope. Suppose
P\_L(f)=P\_L(g), set b=min(f,g), and note P\_L(b)=P\_L(f): the latter
genuine psh function is a minorant of both obstacles, while monotonicity
gives the reverse inequality. Equation (3.2) gives E(b)=E(f)=E(g). For
every real t and every continuous h, monotonicity gives

\begin{verbatim}
E(b+th)<=E(f+th).
\end{verbatim}

The difference has a minimum zero at t=0 and is differentiable there by
(5.1). Hence integral h dmu\_b=integral h dmu\_f.~All original
continuous tests give mu\_b=mu\_f, and the same argument gives
mu\_b=mu\_g. Consequently

\begin{verbatim}
M_L(P_L(f)) := mu_f
\end{verbatim}

is an actual target-independent operator on the genuine
continuous-obstacle envelopes, with raw mass vol(L). It is not defined
by choosing an inverse target, and its model stages are the positive
central coefficients of own143.

One also obtains exact contact orthogonality. The gap f-P\_L(f) is lower
semicontinuous and nonnegative. Take a continuous h\textgreater=0
supported in the open set where this gap exceeds delta\textgreater0. For
0\textless=t\textless=delta/\textbar\textbar h\textbar\textbar,
P\_L(f-th)=P\_L(f), so (3.2) makes E(f-th) constant. Equation (5.1)
implies integral h dmu\_f=0. Compact Hausdorff normality and Radon
regularity imply that mu\_f gives zero mass to \{P\_L(f)\textless f\}.
Taking countably many positive rational delta exhausts this open
strict-gap set. Thus the measure is supported on the contact set in this
precise measure sense, including the reference pole convention.

\subsubsection{6. Exact remaining scope}\label{exact-remaining-scope}

This argument establishes a proposed operator and energy on all genuine
continuous-obstacle envelopes for any actual big rational L in any
dimension. Its all-continuous obstacle derivatives, model invariance,
envelope-fiber descent, positivity, correct raw volume mass, and contact
orthogonality have separate proofs above.

It remains to prove the entire relative E1 construction with mixed and
weighted products, arbitrary monotone nets, weak energy upper
semicontinuity, intrinsic dual targets, actual attainment and all-point
uniqueness. In particular a general normalized psh function or a
bounded-relative psh function has not been proved to be one
continuous-obstacle envelope, and it cannot be silently replaced by an
abstract completion. The original P16.5 and all nonbig extensions remain
unresolved. The imported algebraic and general-envelope dependencies
have no new independent verdict here, and no active native reviewer has
received this argument.

\section{IV. Entire genuine relative energy and weak
compactness}\label{iv.-entire-genuine-relative-energy-and-weak-compactness}

This is an unreviewed continuation of own143--144, with all of their
external algebraic and general-envelope dependencies retained. It
defines energy on actual functions of the entire original psh cone,
proves exact relative-floor characterization, decreasing-net energy
continuity, weak upper semicontinuity and normalized compact energy
sublevels. It does not yet extend the Monge--Ampere operator to that
entire space, prove its mixed/weighted continuity, or solve the original
equation. Original P16.5 remains open.

\subsubsection{1. Genuine operations without an abstract net-flattening
lemma}\label{genuine-operations-without-an-abstract-net-flattening-lemma}

Keep smooth integral projective X over algebraically closed trivially
valued characteristic-zero k, of any cardinality, and any actual big
rational L, theta=c1(L), V=vol(L)\textgreater0. Genuine PSH(theta) has
exactly the original decreasing generically finite FS-net definition,
and V\_theta=P\_theta(0) is the genuine canonical minimal reference
established by the published envelope input of own141. The weak topology
is the original pointwise topology on X\^{}div, including its generic
valuation.

Here are direct constructions of the operations needed below. An actual
ample rational bundle B has a finite generating set in a divisible
power. Its weight-zero FS function is identically zero: at each analytic
point at least one generator is nonzero at its center and has trivial
norm one. Tensoring by these generators therefore embeds a generically
finite FS function in the class of an actual bundle A into that of A+B
with exactly the same function.

Let u,v in PSH(theta) be given by decreasing genuine FS nets phi\_i on
actual L\_i and psi\_j on actual M\_j, with numerical classes tending to
theta. Choose an ample H. For each sufficiently late pair i,j choose
rational epsilon\_i,j\textgreater0 tending to zero such that
L\_i+epsilon\_i,j H-M\_j is ample. Set B\_i,j=L\_i+epsilon\_i,j H. Both
B\_i,j-L\_i and B\_i,j-M\_j are ample. The zero-FS tensor construction
places phi\_i and psi\_j in the same actual B\_i,j without changing
their values. Their maximum is a GF FS function on B\_i,j; its
product-index net is decreasing to max(u,v), and
c1(B\_i,j)-\textgreater theta. This is one actual defining net, so
max(u,v) is genuine. No monotonicity of the numerical epsilon values is
required, since the function values are independent of them.

Real convex combinations are genuine by another single-net construction.
Take a tail on which phi\_i and psi\_j are \textless=C for a fixed
finite C. For 0\textless a\textless1 choose increasing positive rational
a\_l-\textgreater a and b\_l-\textgreater1-a. The actual GF FS functions

\begin{verbatim}
a_l(phi_i-C)+b_l(psi_j-C)+C
\end{verbatim}

are defined by rational tensor powers on actual bundles a\_l L\_i+b\_l
M\_j. They decrease in the three product indices because the shifted
functions are nonpositive. Their numerical classes tend to theta, and
their pointwise limit is au+(1-a)v, including the -infinity convention.
Addition of a real constant is immediate at each GF stage. The cases
a=0,1 are immediate. This proves the needed cone operations directly,
including real coefficients.

Choose an actual ample Omega with Omega-theta ample numerically. On a
tail, Omega-L\_i is ample, and the same zero-FS tensor construction
makes the entire defining net for u a GF FS net on Omega. Thus every
genuine u belongs to PSH(Omega) as the same all-point function. This
does not identify its intrinsic finite energy with ambient ample finite
energy. It does give the standard ample finite-coordinate bound on every
divisorial valuation: for normalized u,
-T\_Omega(v)\textless=u(v)\textless=0. The finite growth of divisorial
valuations follows by the intersection estimate of the complete local
primary Proposition1.28, literal LF922--931. Alternatively the granted
actual-AMPLE G0 supplies this supporting bound within its exact ample
domain. The full actual G0 body remains the sole equation grant; no
nonample equation is granted.

In fact the finite-coordinate bound has a direct stage proof here: every
GF FS function phi\_i on the same actual Omega satisfies
phi\_i(v)\textgreater=phi\_i(v\_triv)-T\_Omega(v), since each nonzero
section in its maximum has m\^{}(-1)v(s)\textless=T\_Omega(v). The
finite maximum inherits this bound. Let the original decreasing net tend
to u; its generic values tend to u(v\_triv), and the bound follows
without an abstract flattening or a new ample regularization step.

\subsubsection{2. Actual continuous-obstacle envelopes decrease to every
genuine
function}\label{actual-continuous-obstacle-envelopes-decrease-to-every-genuine-function}

For any genuine u, let C\_u be the directed set of original continuous
real functions f\textgreater=u, ordered so that functions decrease. It
is nonempty because a genuine u is bounded above, and is directed by
minimum. On the compact Hausdorff X\^{}an an upper semicontinuous
function is the pointwise infimum of its continuous majorants, including
value -infinity. For completeness, at a point x and any real
a\textgreater u(x), choose a neighborhood where u\textless a, a global
upper bound B, and a continuous cutoff that is zero at x and one outside
that neighborhood. The function interpolating from a to max(a,B) is a
majorant with value a at x. Normality of compact Hausdorff spaces
supplies this cutoff.

Since u is a genuine psh minorant of each f,

\begin{verbatim}
u <= P_theta(f) <= f.
\end{verbatim}

Both outer bounds have pointwise infimum u. Therefore the genuine
continuous-obstacle envelopes P\_theta(f) decrease on every original
analytic point to the actual original u. This constructs an
approximation on the original carrier for all genuine u; it does not
replace them by formal completion elements.

Rational model majorants also suffice. Given a continuous majorant f and
epsilon\textgreater0, uniform model density yields a rational model g
with \textbar\textbar g-f\textbar\textbar\textless epsilon/2; add a
rational constant in {[}epsilon/2,epsilon{]} to obtain
g'\textgreater=f.~Rational model majorants are directed by their
model-function minimum and have the same pointwise infimum u. Arbitrary
cardinality and the original net are retained.

\subsubsection{3. Energy of the entire genuine
cone}\label{energy-of-the-entire-genuine-cone}

Use E\_theta\^{}obs on continuous obstacles from own144, with raw mass V
and E\_theta\^{}obs(0)=0. Define on actual genuine functions

\begin{verbatim}
E_theta(u)=inf {E_theta^obs(f): f continuous and f>=u}.    (3.1)
\end{verbatim}

The value lies in R union \{-infinity\}; a constant upper majorant gives
a finite upper bound. It is unchanged if the infimum is restricted to
rational model majorants, by the V-Lipschitz obstacle energy. It is
monotone in u, and shifting the majorants gives

\begin{verbatim}
E_theta(u+c)=E_theta(u)+cV.
\end{verbatim}

It agrees exactly with own144 on every genuine continuous-obstacle
envelope w=P\_theta(g). Indeed any f\textgreater=w has
P\_theta(f)\textgreater=w, and envelope-order monotonicity of own144
gives E\_theta\textsuperscript{obs(f)\textgreater=E\_theta}obs(g).
Conversely g is a continuous majorant of w. Taking the infimum gives
E\_theta(w)=E\_theta\^{}obs(g).

In particular E\_theta(V\_theta)=0 and E\_theta(V\_theta-k)=-kV. For
normalized genuine u, sup u=u(v\_triv)=0 and u\textless=V\_theta, so
E\_theta(u)\textless=0. The normalization of V\_theta is also zero: it
is a minorant of zero and has generic value zero by the published
reference construction.

There is no subtraction of -infinity values in (3.1), and no arbitrary
target measure occurs in the definition.

\subsubsection{4. Exact energy continuity for decreasing
nets}\label{exact-energy-continuity-for-decreasing-nets}

Suppose u\_i is any decreasing net of genuine functions with pointwise
genuine limit u. Write ell=inf\_i E\_theta(u\_i), so
ell\textgreater=E\_theta(u) by monotonicity. Fix a continuous
f\textgreater=u and epsilon\textgreater0. The closed compact sets

\begin{verbatim}
A_i={x: u_i(x)>=f(x)+epsilon}
\end{verbatim}

are decreasing and have empty intersection, because
u\textless=f.~Compactness and directedness imply A\_i is empty for one
index and hence every later index. Thus u\_i\textless=f+epsilon
eventually, so

\begin{verbatim}
ell <= E_theta^obs(f+epsilon)=E_theta^obs(f)+epsilon V.
\end{verbatim}

Take the infimum over f and then epsilon down to zero. This proves

\begin{verbatim}
E_theta(u_i) decreases to E_theta(u).                   (4.1)
\end{verbatim}

If E\_theta(u)=-infinity, apply the same argument to majorants with
arbitrarily negative finite obstacle energy. This is a topological
compact-set argument, not Borel monotone convergence applied to a
possibly uncountable net.

For normalized u set

\begin{verbatim}
u_k=max(u,V_theta-k).
\end{verbatim}

The direct maximum construction makes u\_k genuine, and u\_k decreases
pointwise to u. Where V\_theta=-infinity, u also equals -infinity and
each floor has that same value; no extended subtraction is taken.
Monotonicity gives

\begin{verbatim}
-kV <= E_theta(u_k) <=0,
E_theta(u)=inf_k E_theta(u_k).                           (4.2)
\end{verbatim}

Consequently the actual intrinsic relative finite-energy class

\begin{verbatim}
E_theta^1={u in PSH(theta): E_theta(u)>-infinity}
\end{verbatim}

is exactly the class with uniformly bounded relative-floor energy cost
after generic normalization. The reference singularity type is built
into these genuine floors. This is the entire original cone, not an
ample ambient E1 or a fixed set of continuous envelopes.

\subsubsection{5. Finite-divisorial supporting planes and weak upper
semicontinuity}\label{finite-divisorial-supporting-planes-and-weak-upper-semicontinuity}

For a rational model f\_D let mu\_D be the actual finite divisorial
measure of own143--144, with mass V. Concavity and the derivative of the
continuous obstacle energy yield, for every continuous g,

\begin{verbatim}
E_theta^obs(g) <= E_theta^obs(f_D)+ integral(g-f_D) dmu_D.
\end{verbatim}

Take the full decreasing net of continuous majorants g of a genuine u.
On the finitely many atoms of mu\_D the integrals tend to the finite sum
of u values. The left side tends to (3.1). Hence

\begin{verbatim}
E_theta(u) <= A_D(u),
A_D(u)=E_theta^obs(f_D)+ integral(u-f_D) dmu_D.            (5.1)
\end{verbatim}

Every A\_D(u) is finite: a genuine u is finite on each divisorial or
generic atom, as proved in Section1. If f\_D itself majorizes u, the
integral in A\_D(u) is nonpositive and
A\_D(u)\textless=E\_theta\^{}obs(f\_D). Model majorants suffice in
(3.1), so

\begin{verbatim}
E_theta(u)=inf over all rational models D of A_D(u).      (5.2)
\end{verbatim}

The index set in (5.2) is independent of u. Each A\_D is continuous for
the original pointwise-divisorial topology, since it is a fixed finite
coordinate sum. Thus E\_theta is weakly upper semicontinuous on the
entire genuine cone. The same formula makes it concave on this convex
cone: an infimum of fixed affine functionals is concave. No continuity
of evaluations at nondivisorial poles is assumed.

The general envelope primary input also gives compactness of normalized
PSH(theta) in this exact weak topology, as recorded in own141. Therefore
for every finite B the normalized cost sublevel

\begin{verbatim}
{u in PSH(theta): u(v_triv)=0, E_theta(u)>=-B}
\end{verbatim}

is weakly closed in that compact space, hence compact, and all its
members are intrinsic E1. No metrizability, countability, support
reduction, or ambient finite-energy assumption enters this proof.

\subsubsection{6. Further conclusions not established
here}\label{further-conclusions-not-established-here}

Sections2--5 solve several energy/compactness gates left by own144 on
actual genuine functions. They do not prove that the derivative measures
of different continuous-obstacle majorant nets have the same limit at
every intrinsic E1 endpoint. Uniform derivative control in own144
depends on a uniform bound on the continuous obstacles, which has not
been proved for these majorant nets; bounded relative floors need not
have uniformly bounded absolute continuous obstacles when V\_theta is
singular.

The general-big operator at every E1 function, model-compatible
mixed/weighted products, comparison, finite-dual target
characterization, coercivity and target integral continuity, actual
attainment, and all-point uniqueness remain necessary original
obligations. Concavity alone does not imply the quadratic energy scaling
or the quantitative I estimates used in the nef proofs. They cannot be
imported by replacing zero with V\_theta without proof. The raw
volume-zero cases also remain separately unresolved beyond the earlier
scoped results. Every supporting proposal here remains independently
unreviewed.

\section{V. Truncated volume, the entire operator and outer
derivatives}\label{v.-truncated-volume-the-entire-operator-and-outer-derivatives}

This unreviewed continuation uses the exact section lattices of own143,
obstacle energy and measures of own144, and entire genuine energy of
own145. It proposes the missing target-independent operator on the
entire original relative E1 space, with all original continuous outer
derivatives and weak continuity along decreasing nets and
weak-plus-energy convergence. The proof uses truncated algebraic volumes
to replace the uniform absolute-obstacle bound; it retains all explicit
imported algebraic, envelope-compactness and exact ample G0 supporting
inputs. Mixed products, weighted convergence, pole-nullness, comparison
and the all-target inverse remain open. Original P16.5 remains open.

\subsubsection{1. Exact domain and the truncation
identity}\label{exact-domain-and-the-truncation-identity}

Keep smooth integral projective X of dimension n over any algebraically
closed trivially valued characteristic-zero k, actual big rational L,
V=vol(L)\textgreater0 and d=n+1. Fix a smooth equivariant snc model Y
over X times P1, with F=sum b\_i E\_i, vertical rational D and original
model obstacle f\_D. Put

\begin{verbatim}
alpha_D(c)=rho^*L+D+cF,
g_D(c)=vol_Y(alpha_D(c))/d,
E(D)=E_theta^obs(f_D).
\end{verbatim}

Let w\_m,j be the actual lattice weights of own143 and
r\_m=h\^{}0(X,mL). For every nonnegative rational c the exact
direct-image splitting gives

\begin{verbatim}
h^0(Y,m alpha_D(c))=sum_j max(w_m,j+mc+1,0).
\end{verbatim}

Thus

\begin{verbatim}
g_D(c)=lim_m (n!/m^(n+1))sum_j (w_m,j+mc)_+.           (1.1)
\end{verbatim}

The discarded +1 has error bounded by n!r\_m/m\^{}(n+1), tending to
zero. The required section-volume limit is the stated imported algebraic
premise when alpha\_D(c) is big. If it is not big, its volume is zero,
so the nonnegative normalized section count has limsup zero and hence
limit zero. The integer c-degree conventions are absorbed in the common
divisible subsequence. Continuity extends (1.1)'s consequences to real
c.~On a fixed real vertical-divisor space the same identities extend by
monotone rational coefficient sandwiches and continuity.

Suppose f\_D\textless=0 and E(D)\textgreater=-B for a fixed
B\textgreater=0. Effectivity of model functions means D\textless=0. All
lattice weights are \textless=0. The actual full-weight limit is E(D).
Consequently g\_D is convex and V-Lipschitz as a function of
c\textgreater=0, g\_D(0)=0, and

\begin{verbatim}
g_D(c)>=cV+E(D)>=cV-B.                               (1.2)
\end{verbatim}

Whenever c\textgreater B/V, alpha\_D(c) is big and
vol(alpha\_D(c))\textgreater=d(cV-B)\textgreater0. No lower coefficient
bound for D is assumed. The zero value at c=0 also follows from
alpha\_D(0)\textless=rho\^{}*L, whose volume in dimension n+1 is zero.

\subsubsection{2. A deficient positive model measure and a sharp tail
bound}\label{a-deficient-positive-model-measure-and-a-sharp-tail-bound}

For c\textgreater B/V use the algebraic C1 volume theorem to define

\begin{verbatim}
mu_D,c=sum_i b_i <alpha_D(c)^n>.E_i delta_x_i,
\end{verbatim}

where x\_i is the actual normalized central valuation. Each coefficient
is nonnegative by volume monotonicity under adding the effective E\_i.
Its mass is

\begin{verbatim}
mu_D,c(X^an)=g_D'(c)=<alpha_D(c)^n>.F<=V.              (2.1)
\end{verbatim}

Convexity and (1.2) give

\begin{verbatim}
g_D'(c)>=(g_D(c)-g_D(0))/c>=V-B/c.                   (2.2)
\end{verbatim}

These deficient measures are invariant under model pullback, by the same
exceptional-divisor section/volume derivative proof as own143, now at
this fixed big alpha\_D(c).

More is needed than a bound on the total mass. The actual full-mass
model measure mu\_D of own143 satisfies the coefficientwise measure
inequality

\begin{verbatim}
0<=mu_D,c<=mu_D.                                     (2.3)
\end{verbatim}

To prove it, on the fixed model set

\begin{verbatim}
T_D(c)=E(D)+cV-g_D(c)
      =lim_m (n!/m^(n+1))sum_j min(w_m,j+mc,0).        (2.4)
\end{verbatim}

For a coefficientwise effective vertical increment H, inclusion of the
actual lattices orders every flag weight. The function min(w+mc,0) is
increasing in w. Hence T\_(D+H)(c)\textgreater=T\_D(c). This
monotonicity is valid even though T is usually negative. When
alpha\_D(c) is big, both terms are differentiable under a sufficiently
small increment of any E\_i: E uses its eventual affine-volume
derivative from own143, and g uses the C1 volume differential at this
fixed c.~The right derivative of (2.4) in E\_i is
q\_i(D)-\textless alpha\_D(c)\^{}n\textgreater.E\_i\textgreater=0.
Multiplying by b\_i gives (2.3). This argument does not pair an
arbitrary pseudoeffective curve difference with the nonnef divisor E\_i.

Combining (2.1)--(2.3),

\begin{verbatim}
||mu_D-mu_D,c||_TV
   =V-mu_D,c(X^an)<=B/c.                            (2.5)
\end{verbatim}

Here TV means the norm of the signed measure, not twice the positive
mass. This is an original-carrier tail estimate for the actual proposed
model operator. It is not normalization of the deficient measure.

\subsubsection{3. Uniform big-class margins from volume and an upper
class}\label{uniform-big-class-margins-from-volume-and-an-upper-class}

Fix a rational model direction h=f\_H carried on one fixed smooth model
Z, and a fixed c\textgreater B/V. On any common refinement carrying D
with f\_D\textless=0 and E(D)\textgreater=-B, put

\begin{verbatim}
gamma=rho_Z^*L+cF,
alpha=alpha_D(c),
kappa=d(cV-B)>0.
\end{verbatim}

Then alpha\textless=gamma in pseudoeffective order, gamma is a fixed big
Cartier class on Z, and vol(alpha)\textgreater=kappa. Let T be a fixed
ample Cartier class on Z, pulled to each refinement, and write
Q\_gamma=\textless gamma\^{}n\textgreater. Put M=d
Q\_gamma.T\textgreater=0.

There is a uniform r\_0\textgreater0 such that alpha-r\_0 T is big for
every such D; for instance take r\_0=kappa/{[}2(M+1){]}. To check this,
follow the segment alpha-rT until its first exit from the big cone on
the particular refinement. Inside that cone, positive-product
monotonicity and the nef pairing with T give

\begin{verbatim}
0<=d< (alpha-rT)^n>.T<=M.
\end{verbatim}

The volume differential integrated along the segment gives
vol(alpha-rT)\textgreater=kappa-Mr.~If the segment left the big cone
before r\_0, continuity of volume would make its endpoint volume zero,
contradicting this strictly positive lower bound. Thus it remains big at
r\_0. This argument is uniform across all higher models because gamma
and T are pulled back from the one fixed Z and positive products have
their birational Cartier/Weil pairings.

In particular alpha\textgreater=r\_0 T. Choose a\textgreater0 with
aT\textgreater=+/-H, possible since T is ample on the fixed Z. Then
(a/r\_0)alpha\textgreater=+/-H. For
epsilon=(a/r\_0)\textbar t\textbar\textless1,

\begin{verbatim}
(1-epsilon)alpha<=alpha+tH<=(1+epsilon)alpha.
\end{verbatim}

Positive-product monotonicity and homogeneity now give uniform bounds
for \textless(alpha+tH)\^{}n\textgreater{} between the corresponding
multiples of Q\_alpha. Decompose H=T\_1-T\_2 with fixed nef Cartier
T\_1,T\_2 on Z. Pairing these bounds with T\_1,T\_2 and using
Q\_alpha\textless=Q\_gamma yields

\begin{verbatim}
|[<(alpha+tH)^n>-Q_alpha].H|<=K_c |t|,                (3.1)
\end{verbatim}

for uniformly small t, with K\_c independent of the model and of D. No
nefness of the trace of Q\_alpha is assumed. The fixed-c volume
therefore has the uniform Taylor estimate

\begin{verbatim}
|g_(D+tH)(c)-g_D(c)-t integral h dmu_D,c|
      <=K_c t^2/2.                                  (3.2)
\end{verbatim}

Unlike own144, the hypothesis is an energy bound with an upper obstacle
bound. The obstacles may have arbitrarily negative coefficients and
arbitrarily large absolute norm.

\subsubsection{4. Uniform derivative errors for the full model
energy}\label{uniform-derivative-errors-for-the-full-model-energy}

Let H\_0=\textbar\textbar h\textbar\textbar{}\emph{infinity and take
\textbar r\textbar\textless=delta\textless=1 small enough for (3.1). The
obstacle f}(D+rH) is \textless=\textbar r\textbar H\_0. Subtract
s=\textbar r\textbar H\_0 as a constant to get a nonpositive obstacle
D'=D+rH-sF. Obstacle monotonicity, its V-Lipschitz bound and the
constant identity give

\begin{verbatim}
E(D')>=-B-2V H_0.
\end{verbatim}

The constant shift leaves the full model measure unchanged, while
alpha\_(D+rH)(c)=alpha\_D'(c+s). Thus (2.5), applied with the fixed
B'=B+2V H\_0, gives

\begin{verbatim}
||mu_(D+rH)-mu_(D+rH),c||_TV<=B'/c.                  (4.1)
\end{verbatim}

Choose c larger than B'/V when needed; this only restricts the following
estimates to arbitrarily large fixed c.~For the initial D the sharper
bound is B/c.

E(D+rH) has its actual derivative integral h dmu\_(D+rH) at every r by
own143, using a sufficiently large eventual affine parameter for that
fixed finite segment. Integrate the derivative difference. Equations
(3.1), (4.1) and (2.5) give

\begin{verbatim}
|E(D+tH)-E(D)-t integral h dmu_D|
   <=K_c t^2/2+|t| H_0 (B+B')/c.                    (4.2)
\end{verbatim}

One may first fix a large c and its permitted small-t interval, then let
t tend to zero, and only afterwards let c tend to infinity. No uniform
bound on K\_c as c tends to infinity is claimed or needed. The linear
error coefficient tends to zero. This ordering of limits is the main
replacement for a uniform absolute-obstacle bound.

\subsubsection{5. Decreasing genuine closure within the exact allowed
inputs}\label{decreasing-genuine-closure-within-the-exact-allowed-inputs}

We need genuine membership of decreasing limits without invoking the
disputed abstract varying-class net-flattening lemma. Suppose a
decreasing net of genuine theta-psh functions has finite limiting
generic value and a pointwise limit p.~Embed the whole net as the same
functions in PSH(Omega), for the fixed actual ample Omega of own145, by
its explicit zero-FS construction. The exact genuine decreasing-net
closure and all-point divisorial determination of the sole actual-AMPLE
G0 apply within this ample space. They make p a genuine Omega-psh
function.

On the other hand normalize the net by its generic values. Those values
converge to a finite number. General theta-envelope compactness gives a
cluster point in normalized genuine PSH(theta). Its original divisorial
coordinates coincide with those of the normalized p.~Both are genuine
Omega-psh functions, so the exact ample all-point determination
identifies them everywhere. Thus p is genuine theta-psh on the entire
original carrier. This uses the published general-envelope compactness
of own141 and the actual-ample G0 only in their precisely stated
domains; it gives no new nonample equation grant. The G0 full-body path
and SHA are those recorded in own148.

\subsubsection{6. The entire genuine operator and both outer derivative
signs}\label{the-entire-genuine-operator-and-both-outer-derivative-signs}

Let u be any normalized genuine relative-E1 function, so
E\_theta(u)\textgreater-infinity. Direct the rational model majorants
f\textgreater=u with f\textless=0 by decreasing pointwise order. They
are cofinal, their infimum is u at every original analytic point by
own145, and E(f)\textgreater=E\_theta(u). Put B=-E\_theta(u), so (4.2)
applies uniformly to every such majorant, irrespective of its absolute
norm. Their actual measures mu\_f are positive Radon measures of raw
mass V. Compactness of that original measure space gives cluster points.

Fix an original model direction h and a sufficiently small real t. The
genuine envelopes P\_theta(f+th) decrease pointwise. Their generic
values stay finite: u-\textbar t\textbar H\_0 is a genuine minorant and
V\_theta+\textbar t\textbar H\_0 is a genuine upper bound. Section5
gives a genuine theta-psh pointwise limit. It is exactly P\_theta(u+th):
the limit lies below u+th and is genuine, while every genuine minorant
of u+th lies below every P\_theta(f+th). Therefore

\begin{verbatim}
P_theta(f+th) decreases to the actual P_theta(u+th).
\end{verbatim}

This also proves genuine existence of this usc-obstacle envelope; it
does not assume it to be a continuous-obstacle envelope. Its energy is
finite, since it lies between u-\textbar t\textbar H\_0 and
u+\textbar t\textbar H\_0. Own145 decreasing-net energy continuity gives

\begin{verbatim}
E_theta^obs(f+th) -> E_theta(P_theta(u+th)),
E_theta^obs(f) -> E_theta(u).
\end{verbatim}

For any cluster point mu of mu\_f pass (4.2) through these limits:

\begin{verbatim}
|E_theta(P_theta(u+th))-E_theta(u)-t integral h dmu|
   <=K_c t^2/2+|t| H_0(B+B')/c.                     (6.1)
\end{verbatim}

First let positive and negative t tend to zero for fixed c.~Then let c
tend to infinity. The exact two-sided outer derivative exists and equals
integral h dmu. Every cluster point has this same integral for every
original model direction. Uniform model density in all original C0 and
the common mass V imply that the cluster measure is unique and the whole
majorant net converges weakly.

Define M\_theta(u) to be this unique original positive Radon measure. It
is independent of the chosen cofinal decreasing model-majorant net,
since (6.1) identifies every cluster point by the same outer
derivatives. It is target independent: only actual section/volume model
measures and majorants of the original function u occur in its
definition. For every original continuous h, uniform approximation by
model directions and the V-Lipschitz bound under constant shifts extend
the two-sided formula to

\begin{verbatim}
d/dt at 0 E_theta(P_theta(u+th))=integral h dM_theta(u). (6.2)
\end{verbatim}

The same proof of (6.1) supplies its uniform model-direction error bound
at every genuine u with u\textless=0 and E\_theta(u)\textgreater=-B.
Constant shifts extend the construction and estimates to unnormalized
functions with a fixed finite upper bound. M\_theta(u+c)=M\_theta(u), by
translating the majorants and the constant-shift energy identity.

\subsubsection{7. Ordered energy gaps and agreement with the obstacle
operator}\label{ordered-energy-gaps-and-agreement-with-the-obstacle-operator}

The derivative error in (6.1) has an additional useful consequence. If
model obstacles f\textless=g are bounded above by a fixed constant and
have energies bounded below, then monotonicity gives
E(f+th)\textless=E(g+th). Use both signs of t in (4.2). With
Delta=E(g)-E(f)\textgreater=0, one obtains, for every fixed model
direction h,

\begin{verbatim}
|integral h d(mu_f-mu_g)|
   <=Delta/|t|+K_c |t|+2 epsilon_c,                  (7.1)
\end{verbatim}

where epsilon\_c tends to zero and the constants depend only on the
common upper and energy bounds and the fixed direction. A common
constant shift makes both obstacles nonpositive. Consequently an ordered
pair whose energy gap tends to zero has measures whose original
continuous integrals tend to the same value: first let the energy gap
tend to zero for fixed c,t, then t tend to zero, then c tend to
infinity.

Now let u=P\_theta(g) with g continuous. Choose rational model obstacles
g\_i\textgreater=g tending uniformly to g, and choose cofinal decreasing
model majorants f\_j of u. On their product directed set set
b\_i,j=min(g\_i,f\_j). These are model majorants of u, bounded above by
one fixed constant. They decrease to u if the g\_i are chosen
decreasing; such upper approximations can be arranged by taking finite
minima of uniform upper approximants. Their measures converge to
M\_theta(u) by Section6. Meanwhile

\begin{verbatim}
0<=E(g_i)-E(b_i,j)<=E(g_i)-E_theta(u)->0
\end{verbatim}

uniformly in j. Equation (7.1) compares mu\_g\_i with mu\_b\_i,j. Own144
uniform obstacle convergence gives mu\_g\_i -\textgreater{} the original
obstacle measure mu\_g. The ordered energy-gap bound therefore gives
M\_theta(u)=mu\_g. Thus the entire genuine operator extends the
previously defined continuous-obstacle envelope operator and its actual
model-volume stages. Mere existence of an unrelated supergradient would
not prove this agreement; the ordered estimate does.

\subsubsection{8. Monotone nets and weak-plus-energy
continuity}\label{monotone-nets-and-weak-plus-energy-continuity}

The same ordered estimate holds for genuine u\textless=v in E1 with
common upper and lower energy bounds, by using (6.1) and
P\_theta(u+th)\textless=P\_theta(v+th). Therefore ordered genuine energy
gaps tending to zero force the corresponding measures to converge
against every original continuous observable.

For any decreasing genuine E1 net u\_i with pointwise E1 limit u and
finite generic limit, choose a tail with a common finite upper bound and
subtract that constant. Own145 gives
E\_theta(u\_i)-\textgreater E\_theta(u). The ordered estimate with
u\textless=u\_i gives M\_theta(u\_i)-\textgreater M\_theta(u) weakly on
the full original carrier. No cardinality or sequential assumption is
used.

More generally suppose u\_i -\textgreater{} u in the original divisorial
topology and
E\_theta(u\_i)-\textgreater E\_theta(u)\textgreater-infinity. Put
w\_i=max(u\_i,u), genuine by own145. Its divisorial coordinates converge
to u, and weak energy upper semicontinuity and monotonicity give
E\_theta(w\_i)-\textgreater E\_theta(u). The generic values have a
common finite bound. Thus E\_theta(w\_i)-E\_theta(u\_i) and
E\_theta(w\_i)-E\_theta(u) tend to zero. Apply the ordered estimate to
both pairs and conclude M\_theta(u\_i)-\textgreater M\_theta(u) weakly,
again against all original C0 functions. This is weak-plus-energy
continuity; no unproved equivalence with an I or rooftop metric is
asserted.

\subsubsection{9. Exact unresolved
obligations}\label{exact-unresolved-obligations}

The proposed construction has genuine original functions throughout,
correct raw mass vol(L), target independence, compatibility with the
section-lattice obstacle operator, all-original-C0 two-sided outer
derivatives, and the stated monotone/weak-plus-energy continuity. Its
pivotal new inputs are (2.3), a coefficientwise tail inequality proved
through actual flag weights, and Section3, a uniform big-class margin
obtained from a volume lower bound and a fixed upper class.

The whole dependency chain has not been independently accepted. In
particular the algebraic positive-product order and nef pairings, the
true section-volume asymptotics, the general numerical envelope
compactness, and the exact supporting applications of ample G0 must
retain their full checked hypotheses. Nothing here proves relative
pole-nullness, weighted finite-energy integrability, positive mixed
products, comparison, domination, quantitative coercivity of the full
intrinsic dual target class, actual attainment for arbitrary such
targets, or general all-point uniqueness. These remain required to solve
original P16.5. The active native124/125 and frozen126/127 inputs have
not been changed or given this argument as hot feedback.

\section{VI. Complete independent reconstruction and final
scope}\label{vi.-complete-independent-reconstruction-and-final-scope}

Scoped PASS. The chain establishes the requested energy and
target-independent operator for every smooth integral projective
original X of positive dimension, every algebraically closed trivially
valued characteristic-zero k of arbitrary cardinality, and every actual
big rational L. Nefness, finite generation, and a Zariski decomposition
are unnecessary. This conclusion leaves original P16.5 open.

The external premises are indexed below. Passage −1 records the sole
task grant and is deliberately excluded from the external-premise array.
I use G0 only for genuine psh closure and divisorial determination
inside an actual ample class; its ample energy, intersection mass, mixed
products, and inverse theorem are never transferred to L.

\begin{enumerate}
\def\labelenumi{\arabic{enumi}.}
\item
  General envelopes and the original reference. I opened the
  \href{https://arxiv.org/pdf/2206.07183v2}{BJ Addendum}, including
  Lemmas 1.3--1.7, the characteristic-zero proof of Theorem A, and
  Corollaries 2.1--2.2. In the needed big rational case, the
  multiplier-ideal approximants are already one actual defining net on
  L+A/m. Subadditivity and the shift making the auxiliary model function
  nonnegative give decrease in the divisibility order. A shifted nonzero
  section supplies a genuine minorant and prevents generic collapse.
  This membership step requires no abstract varying-class flattening.
  The reference and compactness inputs retain the smooth original domain
  and arbitrary cardinality. In particular V\_L=P\_L(0) is genuine, has
  generic value zero, and dominates every normalized genuine potential.
  For usc obstacles with a genuine minorant, envelope regularity and the
  obstacle inequality make the regularized supremum itself a competitor,
  so the envelope equals the unregularized supremum everywhere.
\item
  Algebraic inputs and their scope. I checked the pertinent
  positive-product construction, continuity, concavity, and
  differentiability units in
  \href{https://arxiv.org/pdf/math/0608260v2}{BFJ}. Every volume
  differential in the argument is applied on a smooth projective variety
  of dimension d=n+1, at a big Cartier class. The order comparisons
  concern Cartier classes and compatible Weil positive products. Their
  intersections with fixed pulled nef Cartier directions are legitimate
  nonnegative pairings; arbitrary nonnef central components are not used
  as order-preserving tests. Full section limits follow from the
  complete-series proof in
  \href{https://www.math.stonybrook.edu/robert.lazarsfeld/Reprints/Laz.Mustata.CBALS.pdf}{Lazarsfeld--Mustaţă},
  §§2.1--2.2. Its uncountable-field convention causes no restriction
  here: extending k to an uncountable algebraically closed field
  preserves each actual section space by flat base change, hence
  preserves the entire numerical sequence and its limit. Clearing a
  rational denominator gives the required divisible-degree limits. No
  finite-generation assumption enters.
\item
  Actual lattices, energy, and multiplicities. Proper direct image is
  coherent and torsion free over the smooth base curve, hence locally
  free. The fixed product trivialization and equivariance turn its
  lattice into a finite flag in R\_m=H\^{}0(X,mL). An adapted basis
  gives M\_m=⊕\_j k{[}t{]}t\^{}(−w\_m,j)s\_m,j, with ma≤w\_m,j≤mb.
  Gluing at infinity therefore gives the exact formula
  h\^{}0(Y,mα\_D(c))=∑\_j max(w\_m,j+mc+1,0). This is the all-c formula.
  The untruncated affine formula applies only once c\textgreater−a.
  Section asymptotics then give vol(α\_D(c))=dcV+dE(D), where
  V=vol\_X(L), and E(D)=lim\_m n!m\^{}(−d)∑\_j w\_m,j. Thus E(0)=0,
  constant shifts add sV, coefficientwise effective increments increase
  energy, and the sup-norm bound is V-Lipschitz. Expanding volume-root
  concavity at large c proves ordinary concavity of E, without asserting
  ordinary volume concavity.
\end{enumerate}

For sufficiently large c, put q\_i(D)=⟨α\_D(c)\^{}n⟩·E\_i. Adding the
effective E\_i increases volume, so its derivative proves q\_i(D)≥0.
Differentiation in F=∑\_i b\_iE\_i gives ∑\_i b\_iq\_i(D)=V. A common
eventual-affine half-line for D+tE\_i proves independence of each q\_i
from c, rather than just independence of their sum. Consequently
μ\_D=∑\emph{i b\_iq\_i(D)δ}(x\_i) is positive and has raw mass V. Since
f\_H(x\_i)=h\_i/b\_i, differentiation gives dE(D+tH)/dt\textbar\_0=∫f\_H
dμ\_D, retaining every central multiplicity.

\begin{enumerate}
\def\labelenumi{\arabic{enumi}.}
\setcounter{enumi}{3}
\item
  Birational invariance of both measures. For a refinement f:Y′→Y and
  effective exceptional G, normality gives f\_\emph{O\_Y′(kG)=O\_Y:
  permitted poles lie only above codimension at least two, and the
  height-one characterization of regularity excludes them downstairs.
  Projection formula therefore preserves sections and volume after
  adding nonnegative exceptional multiples to a pulled class. At a big
  class, differentiation makes every new exceptional coefficient zero.
  For a strict transform E\_i′, the effective exceptional difference
  f\^{}}E\_i−E\_i′ has zero pairing, so its coefficient equals the
  original coefficient. Strict transforms preserve b\_i and x\_i. This
  proves invariance of the full μ\_D and, separately, of μ\_D,c at every
  fixed c where α\_D(c) is big. Common smooth equivariant snc
  refinements exist: close the common generic graph and apply functorial
  resolution to the central pair. The action and projection maps are
  smooth surjective, so functoriality supplies equivariance; resolution
  changes nothing away from the central fiber.
\item
  Section energy and genuine envelope-fiber descent. A single nonzero
  actual section gives a genuine generically finite FS function. Hence
  W\_m(f;s)≥λ is equivalent, on the whole original carrier, to
  q\_s+λ/m≤f and to q\_s+λ/m≤P\_L(f). For model obstacles, this is
  exactly the central regularity condition for st\^{}(−λ). Horizontal
  divisors introduce no poles, and normality upgrades the prime-divisor
  test to regularity. Thus the genuine envelope determines every actual
  filtered coefficient space, in every divisible degree. Ordered
  envelopes order the flag weights without requiring a simultaneously
  adapted basis.
\end{enumerate}

The flag-energy expression e\_m(f)=n!m\^{}(−d)∑\_j floor W\_m(f;s\_m,j)
consequently agrees with model energy and has the stated uniform
comparison bound. Model density and the full section limit establish its
limit E\^{}obs(f) for every original continuous f.~This functional is
finite, concave, monotone, V-Lipschitz, translation equivariant, and
ordered by genuine envelopes. In particular equal envelopes have equal
energy.

Own144's bounded-obstacle Taylor estimate is valid across all
refinements. Its fixed big lower class gives the multiplicative class
sandwich; positive-product order and homogeneity, paired separately with
the two fixed nef summands of H, give the uniform derivative estimate.
Positive fixed-mass measures then have unique weak clusters, determined
by the two-sided derivatives against model functions and hence all
original continuous functions. The minimum-obstacle argument descends
these measures through equal genuine envelopes. The contact argument
also checks out: continuous nonnegative tests supported inside a strict
obstacle gap permit lowering the obstacle without changing its envelope,
and Radon regularity exhausts that gap.

\begin{enumerate}
\def\labelenumi{\arabic{enumi}.}
\setcounter{enumi}{5}
\tightlist
\item
  Entire genuine relative energy. The cone operations in own145 use
  actual defining nets. Tensoring with weight-zero generators changes
  the bundle while preserving the scalar function; common ample
  increments therefore produce genuine maxima. Increasing rational
  coefficients applied to nonpositive shifted stages produce real convex
  combinations. Every genuine potential embeds, as the same all-point
  function, in one fixed actual ample Ω. If finite FS stages are needed
  for the exact ample G0 convention, append the Ω generators at weight
  −k: max(φ\_i,−k) is a single decreasing product-index net of finite FS
  functions. This does not assert ambient ample finite energy.
\end{enumerate}

Continuous majorants, and then rational model majorants, have pointwise
infimum u on every analytic point. Therefore E\_L(u)=inf\_(f continuous,
f≥u) E\^{}obs(f) is defined on the entire actual cone. It agrees with
obstacle energy on P\_L(f), satisfies E\_L(V\_L)=0 and
E\_L(u+c)=E\_L(u)+cV, and never subtracts two infinite values.

For u\_i↓u with genuine endpoint, the compact decreasing sets
\{u\_i≥f+ε\}, for a continuous majorant f of u, eventually become empty.
This proves decreasing-net energy continuity, including energy −∞.
Applied to u\_k=max(u,V\_L−k), it gives E\_L(u)=inf\_k E\_L(u\_k). After
generic normalization, −kV≤E\_L(u\_k)≤0, so finite energy is exactly
uniformly bounded relative-floor energy cost.

Concavity and the model derivatives give supporting planes
A\_D(u)=E\^{}obs(f\_D)+∫(u−f\_D)dμ\_D. Their atoms are divisorial or
generic, and the actual-ample section-growth bound makes every
coordinate finite. Passing continuous majorants down to u gives
E\_L(u)≤A\_D(u); model majorants give the reverse infimum equality
E\_L(u)=inf\_D A\_D(u). Each plane is a finite continuous
divisorial-coordinate sum. Energy is consequently weakly upper
semicontinuous and concave, and normalized finite energy-cost sublevels
are compact in the original weak topology.

\begin{enumerate}
\def\labelenumi{\arabic{enumi}.}
\setcounter{enumi}{6}
\tightlist
\item
  Truncation and the coefficientwise tail. The all-c lattice formula,
  with its +1 error bounded by n!r\_m/m\^{}d→0, gives g\_D(c)=lim\_m
  n!m\^{}(−d)∑\emph{j(w\_m,j+mc)}+. If α\_D(c) is nonbig, volume zero
  makes the nonnegative normalized section count have limsup zero and
  therefore limit zero. Rational c is treated first; real consequences
  follow by continuity.
\end{enumerate}

For D≤0 and E(D)≥−B, all weights are nonpositive. Thus g\_D(0)=0, g\_D
is convex and V-Lipschitz, and cV−B≤g\_D(c)≤cV. For c\textgreater B/V,
α\_D(c) is big. Volume monotonicity under each effective E\_i proves
positivity of μ\_D,c, and differentiation in F gives its mass g\_D′(c).
Convexity gives V−B/c≤g\_D′(c)≤V.

The indispensable stronger inequality is also valid. The identity
T\_D(c)=E(D)+cV−g\_D(c)=lim\_m n!m\^{}(−d)∑\emph{j min(w\_m,j+mc,0)
holds without a sign restriction on an incremented D. Effective lattice
inclusion orders every actual flag weight, so T}(D+H)(c)≥T\_D(c).
Differentiating this monotonicity in E\_i at a big α\_D(c) gives
q\_i(D)−⟨α\_D(c)\^{}n⟩·E\_i≥0. Hence 0≤μ\_D,c≤μ\_D coefficientwise and
\textbar\textbar μ\_D−μ\_D,c\textbar\textbar\_TV≤B/c.~This conclusion
uses ordered weights, not a curve-order pairing with E\_i.

\begin{enumerate}
\def\labelenumi{\arabic{enumi}.}
\setcounter{enumi}{7}
\tightlist
\item
  Uniform margin and the full-energy derivative error. For α=α\_D(c)≤γ
  with vol(α)≥κ=d(cV−B), follow α−rT until its first possible boundary
  point. Inside the big cone, positive-product order and the pulled nef
  T bound the volume loss by M=d⟨γ\^{}n⟩·T. The integrated differential
  gives vol(α−rT)≥κ−Mr.~Continuity at a boundary would force zero
  volume, contradicting this bound before r\_0=κ/{[}2(M+1){]}. Thus
  α−r\_0T remains big on every refinement, with one common r\_0.
\end{enumerate}

Choose aT≥±H on the fixed direction model. Then α≥r\_0T gives
(1−ε)α≤α+tH≤(1+ε)α, ε=(a/r\_0)\textbar t\textbar. Pairing the resulting
positive-product bounds with fixed nef T\_1,T\_2, where H=T\_1−T\_2,
proves a uniform fixed-c Taylor estimate. Perturbed obstacles are
correctly normalized by subtracting \textbar r\textbar{}
\textbar\textbar h\textbar\textbar∞. Their energy is bounded below by
−B′, B′=B+2V\textbar\textbar h\textbar\textbar∞, and their deficient
class is tested at c+\textbar r\textbar{}
\textbar\textbar h\textbar\textbar∞. Constant shifts preserve the full
measure. Integrating derivatives and using both tail bounds yields the
uniform error
K\_ct\^{}2/2+\textbar t\textbar\textbar\textbar h\textbar\textbar∞(B+B′)/c.~The
argument fixes c, sends t→0, and only then sends c→∞; no uniform bound
on K\_c as c grows is required.

\begin{enumerate}
\def\labelenumi{\arabic{enumi}.}
\setcounter{enumi}{8}
\tightlist
\item
  Genuine closure and the entire operator. For a decreasing genuine
  θ-net with finite limiting generic value, ample G0 first makes its
  all-point limit genuine in Ω. Normalize the original net by its
  convergent generic values and use the general θ compactness input to
  obtain a genuine θ cluster. Its divisorial coordinates coincide with
  the ample limit. All-point determination is then applied only to the
  two genuine Ω-potentials, establishing θ membership of the actual
  limit.
\end{enumerate}

For u in the intrinsic relative E1 class, nonpositive model majorants
carry one uniform energy lower bound. For every model direction h, their
perturbed envelopes decrease to a genuine limit. The limit lies below
u+th, while every genuine minorant of u+th lies below each envelope,
proving whole-carrier equality with P\_L(u+th). The bounds by
u±\textbar t\textbar\textbar\textbar h\textbar\textbar∞ keep its energy
finite. Decreasing energy continuity therefore passes the uniform model
error to this actual envelope energy.

Both quotient signs followed by c→∞ give the derivative. Every weak
cluster has the same model-test integral; uniform density and fixed mass
identify a unique positive Radon cluster on all original C0 tests. This
defines M\_L(u), of raw mass vol(L), independently of targets and
majorant choices. Uniform approximation of the direction, with energy
error at most
V\textbar t\textbar\textbar\textbar h−h\_m\textbar\textbar∞, extends the
two-sided formula to every continuous h. The same argument applies to
any decreasing model-majorant net with actual pointwise endpoint u and a
finite upper tail, even if the net is not order-cofinal.

\begin{enumerate}
\def\labelenumi{\arabic{enumi}.}
\setcounter{enumi}{9}
\tightlist
\item
  Agreement and continuity. For ordered pairs, the two derivative signs
  give the bound Δ/\textbar t\textbar+K\_c\textbar t\textbar+2ε\_c on
  the difference of model-test integrals. Its limit order is energy
  gap→0, then t→0, then c→∞. To compare u=P\_L(g) with the earlier
  obstacle operator, take decreasing uniform upper model approximations
  g\_i of the single continuous g and model majorants f\_j of u. The
  model minimum b\_i,j=min(g\_i,f\_j) still majorizes u, and
  0≤E(g\_i)−E(b\_i,j)≤E(g\_i)−E(u)→0 uniformly in j. Thus the ordered
  estimate proves M\_L(u)=μ\_g. The extension is genuinely compatible
  with the prior section/volume measures.
\end{enumerate}

The ordered estimate passes to genuine E1 pairs. Decreasing-net energy
continuity then proves M\_L(u\_i)→M\_L(u) weakly for every decreasing
genuine net with E1 endpoint. A decreasing net with a genuine weak
endpoint has that same all-point endpoint by ample closure and
determination. Finally, if u\_i→u divisorially and E\_L(u\_i)→E\_L(u),
genuine w\_i=max(u\_i,u) converge divisorially to u. Monotonicity and
upper semicontinuity squeeze E\_L(w\_i) to E\_L(u). Both ordered energy
gaps vanish, proving weak-plus-energy operator continuity against every
original continuous observable. No countable replacement is used.

Harmless documentary corrections: own144's references to section `poles'
should read section zeros, where log\textbar s\textbar=−∞. Own150's
unavailable own148 path/hash reference is unnecessary because G0 is
embedded and explicitly granted here. The full-section-limit evidence is
supplied above rather than inferred from BFJ's limsup definition. No
substantial mathematical repair was needed. No mixed, weighted,
pole-nullness, comparison, or all-target inverse conclusion is included
in this PASS.

\section{VII. Documentary corrections and public
evidence}\label{vii.-documentary-corrections-and-public-evidence}

The references in own144 to section poles mean section zeros, where the
scalar is minus infinity. The unavailable own148 path mentioned in
own150 is not a premise: its needed ample statements are in the exact G0
body. The independent report supplies the complete-series section-limit
premise as explained in Section0. These documentary clarifications
change neither the frozen proof nor the accepted theorem.

The canonical native report has 11 external premises and 248 passage
records, many repeating the same indexed statement with a short quote.
Its task-grant passage is the final record at urn:xean:task; its
numerical premise index0 is a serialization association and does not
turn the ample grant into the addendum theorem. The public projection
preserves the report, statements, links, ids and indices, while omitting
all verbatim source quotes. The local raw canonical result is bound by
its exact hash in the ledger. No private transcripts, logs, inputs or
databases accompany this reader.

\end{document}
